\chapter{Measuring FRFP Health: Metrics, Indicators, and Checklists}
\label{ch:frfp-metrics}

This chapter offers a practical toolkit for assessing how well a system
(or organisation) aligns with FRFP principles.
It does \emph{not} define a single scalar ``FRFP score''.
Instead, it gathers:

\begin{itemize}
  \item structural indicators for explicit semantics, boundaries, safe
        horizons, and governance;
  \item simple maturity levels;
  \item concrete checklists you can apply to an existing or proposed
        system.
\end{itemize}

The intent is diagnostic and comparative:
to help you see where a system is strong or weak in the FRFP sense,
and where design or governance changes are most urgent.

\section{Why FRFP Metrics Are Different}

Traditional evaluation often asks:
\emph{How accurate is the model? How fast is the system?}
FRFP asks additional questions:

\begin{itemize}
  \item Is the explicit mapping the model implements well-defined and
        appropriate for the tacit function we care about?
  \item Where do humans or institutions actually endorse decisions, and
        are those boundaries meaningful?
  \item How far can the system run (in time, scale, and domain) before
        someone must re-ground assumptions?
  \item Which parts of governance are de facto automated, and which
        remain genuinely tacit?
\end{itemize}

FRFP ``metrics'' are therefore mostly \emph{structural indicators} and
\emph{maturity assessments}, not just numbers you can compute from logs.
They are intended to be used alongside conventional measures of
performance and safety, not instead of them.

\section{Core FRFP Indicator Families}

We group indicators into five families:

\begin{enumerate}
  \item explicit-task semantics;
  \item boundary design and implementation;
  \item safe horizons and dynamics;
  \item governance and institutional operators;
  \item tacit-skill and norm preservation.
\end{enumerate}

For each family, we provide:

\begin{itemize}
  \item key questions;
  \item example indicators;
  \item rough maturity levels.
\end{itemize}

\subsection{Explicit Semantics Indicators}

\paragraph{Key questions.}

\begin{itemize}
  \item Is there a clear, written description of what the system is
        trying to compute?
  \item Is there a plausible ground truth or evaluation procedure?
  \item Are failure modes understood as \emph{semantic} (wrong
        question) or \emph{statistical} (noise, data limits)?
\end{itemize}

\paragraph{Indicative signals.}

\begin{itemize}
  \item Presence of a concise ``task definition'' document that could be
        given to a sceptical domain expert.
  \item Availability of representative test sets or benchmarks that
        reflect real use, not just training data splits.
  \item Explicit list of assumptions and known \emph{out-of-scope}
        cases.
\end{itemize}

\paragraph{Maturity levels.}

\begin{description}
  \item[Level 0 (unclear):]
    Task definition is vague
    (e.g.\ ``rank the best content'', ``optimise risk-adjusted
    returns''), with no precise mapping or ground truth.
  \item[Level 1 (local clarity):]
    Task is clearly defined for a narrow slice of cases, but the system
    is already used more broadly.
  \item[Level 2 (aligned core):]
    Task is well-defined and most usage stays inside that scope; there
    is a known set of out-of-scope cases.
  \item[Level 3 (strong semantics):]
    Task is sharply defined; there is robust evaluation and continuous
    checking that usage remains within validated semantics.
\end{description}

\subsection{Boundary Indicators}

\paragraph{Key questions.}

\begin{itemize}
  \item Can we name the moments where humans or institutions endorse
        or commit?
  \item Are those moments technically and procedurally real, or merely
        rhetorical?
  \item Does the system log and support those boundary events?
\end{itemize}

\paragraph{Indicative signals.}

\begin{itemize}
  \item Clear documentation of which decisions AI systems may take
        autonomously and which require human sign-off.
  \item User interfaces where boundary actions (approve, reject, refer,
        sign, file, deploy) are explicit and foregrounded.
  \item Logging that ties each boundary event to:
        \begin{itemize}
          \item responsible agent;
          \item information available at decision time;
          \item relevant model outputs and uncertainties.
        \end{itemize}
\end{itemize}

\paragraph{Maturity levels.}

\begin{description}
  \item[Level 0 (implicit):]
    Humans are nominally ``in the loop'', but in practice no clear
    boundary points exist; model outputs are simply followed.
  \item[Level 1 (documented):]
    Policies describe boundaries (e.g.\ ``doctor is responsible''),
    but systems and logs do not enforce or reveal them.
  \item[Level 2 (implemented):]
    Boundary events are explicit in UI and logs; humans do endorse, but
    information quality and incentives may still be weak.
  \item[Level 3 (robust):]
    Boundaries are explicit, instrumented, and aligned with incentives;
    external reviewers can reconstruct who endorsed what and on what
    basis.
\end{description}

\subsection{Safe-Horizon Indicators}

\paragraph{Key questions.}

\begin{itemize}
  \item How long can the system run without being deeply re-examined?
  \item How many people, assets, or decisions can it affect before
        someone is required to reconsider design assumptions?
  \item Are there explicit triggers to pause, review, or redesign?
\end{itemize}

\paragraph{Indicative signals.}

\begin{itemize}
  \item Existence of phase-gated deployment plans:
        pilots, limited roll-outs, full deployment.
  \item Formal review intervals for models and programmes (e.g.\ yearly
        clinical revalidation, quarterly strategy deep dives).
  \item Predefined caps on scale (exposure, user count, jurisdiction)
        before further governance action is required.
\end{itemize}

\paragraph{Maturity levels.}

\begin{description}
  \item[Level 0 (unbounded):]
    Once deployed, the system runs indefinitely and scales freely; no
    planned re-grounding.
  \item[Level 1 (reactive):]
    Reviews occur only after problems or complaints; no explicit
    horizons.
  \item[Level 2 (scheduled):]
    Regular reviews and audits exist, but may not be tied to explicit
    caps or triggers.
  \item[Level 3 (engineered horizons):]
    Safe horizons are explicit; hitting a horizon or trigger forces
    structured re-grounding and possible redesign.
\end{description}

\subsection{Governance and Institutional Indicators}

\paragraph{Key questions.}

\begin{itemize}
  \item Who can change the system’s objectives, boundaries, and
        horizons?
  \item What information do they see?
  \item Are they independent of the teams building and operating the
        system?
\end{itemize}

\paragraph{Indicative signals.}

\begin{itemize}
  \item Existence of a named governance body (committee, board, panel)
        with scope over the system.
  \item Published charters or terms of reference for that body.
  \item Evidence that governance decisions affect:
        \begin{itemize}
          \item what the system is allowed to decide;
          \item where and for whom it may operate;
          \item when it must be paused or redesigned.
        \end{itemize}
  \item Use of case studies and tacit evidence (frontline feedback,
        near-miss reports), not just aggregate metrics.
\end{itemize}

\paragraph{Maturity levels.}

\begin{description}
  \item[Level 0 (absent):]
    No dedicated governance; product or line managers implicitly play
    this role.
  \item[Level 1 (nominal):]
    A committee exists on paper but has limited information or
    authority; decisions remain driven by metrics and local teams.
  \item[Level 2 (active):]
    Governance has real influence; it can halt, modify, or constrain
    deployments, but may lack a structured FRFP vocabulary.
  \item[Level 3 (FRFP-aware):]
    Governance bodies explicitly reason in terms of \(E\), \(T\),
    boundaries, and horizons; they are structurally independent and
    use both metrics and tacit input.
\end{description}

\subsection{Tacit Skill and Norm Preservation Indicators}

\paragraph{Key questions.}

\begin{itemize}
  \item Are human skills and norms being preserved, or eroded, by AI
        use?
  \item Do users retain the ability and habit of questioning the
        system?
  \item Is the organisation tracking these changes?
\end{itemize}

\paragraph{Indicative signals.}

\begin{itemize}
  \item Training and onboarding explicitly address:
        \begin{itemize}
          \item when to trust and when to doubt AI outputs;
          \item how to verify or override them;
          \item examples of failures and near-misses.
        \end{itemize}
  \item Periodic assessment of human skills (e.g.\ sample cases where
        staff must solve tasks without AI support).
  \item Monitoring of override rates and rationales:
        are humans still actually exercising judgement?
\end{itemize}

\paragraph{Maturity levels.}

\begin{description}
  \item[Level 0 (ignored):]
    No attention to tacit skill or norm drift; AI is treated purely as a
    productivity boost.
  \item[Level 1 (acknowledged):]
    Informal awareness and discussion, but no systematic measurement or
    design.
  \item[Level 2 (monitored):]
    Some indicators (override rates, training, skill assessments) are
    tracked and inform local adjustments.
  \item[Level 3 (managed):]
    Tacit preservation is a design goal; systems and workflows are
    explicitly shaped to maintain critical human capabilities over
    time.
\end{description}

\section{An FRFP Maturity Snapshot}

For quick diagnosis, you can summarise a system’s FRFP posture in a
table:

\begin{center}
\begin{tabular}{lccc}
\toprule
\textbf{Dimension} & \textbf{Level 0} & \textbf{Level 1--2} & \textbf{Level 3} \\
\midrule
Explicit semantics     & unclear      & partial / local     & strong, validated \\
Boundaries             & implicit     & documented / UI     & explicit, logged \\
Safe horizons          & none         & ad hoc / scheduled  & engineered        \\
Governance             & absent       & nominal / active    & FRFP-aware        \\
Tacit preservation     & ignored      & acknowledged / mon. & actively managed  \\
\bottomrule
\end{tabular}
\end{center}

The goal is not perfection, but awareness.
Many real systems will be at Level~1 or~2 on most dimensions.
The value of this snapshot is in:

\begin{itemize}
  \item making gaps visible (e.g.\ strong semantics but weak horizons);
  \item helping prioritise interventions;
  \item providing a common language across technical and governance
        stakeholders.
\end{itemize}

\section{Checklists}

We now give three concrete checklists you can use directly:

\begin{enumerate}
  \item a system design checklist (pre-deployment);
  \item an operational monitoring checklist (post-deployment);
  \item an institutional governance checklist.
\end{enumerate}

Each question is essentially Yes/No (or ``not applicable''); the aim is
to trigger discussion and action, not to produce a numeric score.

\subsection{System Design Checklist (Pre-Deployment)}

\begin{enumerate}
  \item \textbf{Task semantics}
        \begin{enumerate}
          \item Have we written down, in one paragraph, what function
                the system is supposed to approximate?
          \item Is there any plausible ground truth or external
                evaluation procedure for that function?
          \item Have we listed the main types of cases that are
                \emph{out of scope}?
        \end{enumerate}
  \item \textbf{Data and evaluation}
        \begin{enumerate}
          \item Do we have test data that reflects real deployment
                conditions, not just training distributions?
          \item Are evaluation metrics aligned with what domain experts
                care about?
        \end{enumerate}
  \item \textbf{Boundaries}
        \begin{enumerate}
          \item Have we identified the specific decisions the AI will
                make vs.\ those reserved for humans?
          \item Are boundary events (approvals, sign-offs, filings,
                orders) explicitly represented in the design?
          \item Will our logs capture who endorsed what, and when?
        \end{enumerate}
  \item \textbf{Uncertainty and abstention}
        \begin{enumerate}
          \item Can the system say ``I don't know'' or defer to humans
                when inputs are low-quality or out-of-domain?
          \item Are those deferrals routed to appropriate tacit
                decision-makers?
        \end{enumerate}
  \item \textbf{Safe horizons}
        \begin{enumerate}
          \item Have we defined initial limits on:
                \begin{itemize}
                  \item time (how long the system can run before a deep
                        review)?
                  \item scale (how many users, patients, trades, etc.)?
                  \item domain (which segments or contexts)?
                \end{itemize}
          \item Are there explicit triggers for pausing, rolling back,
                or redesigning (performance drops, incident rates,
                context changes)?
        \end{enumerate}
  \item \textbf{Governance}
        \begin{enumerate}
          \item Is there a named body or role with authority over this
                system’s deployment and evolution?
          \item Have we agreed what information they will see and how
                often they will act?
        \end{enumerate}
  \item \textbf{Tacit preservation}
        \begin{enumerate}
          \item Have we identified which human skills and norms are
                critical to preserve (e.g.\ verification, domain
                intuition)?
          \item Does the design explicitly support those skills
                (e.g.\ by requiring review, offering raw data, or
                including training)?
        \end{enumerate}
\end{enumerate}

\subsection{Operational Monitoring Checklist (Post-Deployment)}

\begin{enumerate}
  \item \textbf{Usage}
        \begin{enumerate}
          \item Is actual use consistent with the intended task and
                domain?
          \item Are there emerging use cases we did not anticipate?
        \end{enumerate}
  \item \textbf{Performance and drift}
        \begin{enumerate}
          \item Are key performance metrics stable in the deployed
                setting?
          \item Have we checked for distribution shift in inputs or
                outcomes?
        \end{enumerate}
  \item \textbf{Boundaries in practice}
        \begin{enumerate}
          \item Are humans still making the decisions they are supposed
                to make, or has the AI effectively taken over?
          \item Do override rates and patterns suggest genuine human
                judgement, or rubber-stamping?
        \end{enumerate}
  \item \textbf{Horizon adherence}
        \begin{enumerate}
          \item Are we still inside the originally defined safe horizons
                (time, scale, domain)?
          \item If not, has any explicit decision been taken to extend
                them?
        \end{enumerate}
  \item \textbf{Incidents and near-misses}
        \begin{enumerate}
          \item Are incidents and near-misses being logged and
                investigated?
          \item Do we see repeated patterns that suggest deeper design
                issues?
        \end{enumerate}
  \item \textbf{Human feedback}
        \begin{enumerate}
          \item Are frontline users reporting concerns or workarounds?
          \item Have we incorporated their tacit knowledge back into
                design updates?
        \end{enumerate}
  \item \textbf{External context}
        \begin{enumerate}
          \item Have regulations, norms, or competitive conditions
                changed in ways that affect admissible runs?
          \item Have we adjusted the system or governance in response?
        \end{enumerate}
\end{enumerate}

\subsection{Institutional Governance Checklist}

\begin{enumerate}
  \item \textbf{Mandate}
        \begin{enumerate}
          \item Does a governance body have an explicit mandate over the
                system (scope, powers, responsibilities)?
          \item Is that mandate documented and communicated to
                stakeholders?
        \end{enumerate}
  \item \textbf{Information}
        \begin{enumerate}
          \item Does the body see both quantitative metrics and
                qualitative case reports?
          \item Are FRFP-relevant artefacts (boundaries, horizons,
                incident logs) part of its regular agenda?
        \end{enumerate}
  \item \textbf{Authority}
        \begin{enumerate}
          \item Can the body actually change:
                \begin{itemize}
                  \item where and how the system is used?
                  \item what tasks it may perform?
                  \item when it must be paused or withdrawn?
                \end{itemize}
          \item Are its decisions binding on operational teams?
        \end{enumerate}
  \item \textbf{Independence and composition}
        \begin{enumerate}
          \item Is membership cross-functional
                (technical, domain experts, legal, ethics, affected
                stakeholders)?
          \item Is it sufficiently independent from the teams whose
                incentives are tied to system growth or metrics?
        \end{enumerate}
  \item \textbf{Processes}
        \begin{enumerate}
          \item Are there scheduled reviews (horizon checkpoints) as
                well as incident-driven meetings?
          \item Are decisions minuted, with rationale and links to
                specific evidence?
        \end{enumerate}
  \item \textbf{External linkage}
        \begin{enumerate}
          \item Does the body interact with external operators
                (regulators, professional societies, public bodies) as
                needed?
          \item Are external developments (new standards, regulations,
                prominent failures elsewhere) fed back into governance?
        \end{enumerate}
\end{enumerate}

\section{How to Use This Chapter}

In practice, the recommended workflow is:

\begin{enumerate}
  \item Pick a concrete system (existing or proposed).
  \item Use the \emph{System Design Checklist} to assess its FRFP
        structure.
  \item Assign rough maturity levels on each dimension
        (semantics, boundaries, horizons, governance, tacit
        preservation).
  \item Identify 2--3 most concerning gaps.
  \item Consult the pattern and anti-pattern chapter to choose design
        moves that address those gaps.
  \item Revisit periodically with the \emph{Operational Monitoring} and
        \emph{Governance} checklists.
\end{enumerate}

The metrics and indicators here are deliberately simple.
They are meant to be used quickly in real organisations, as conversation
starters and design aids, not as precise scientific instruments.
Their value lies in making FRFP questions unavoidable:
for any serious system, someone should be able to say, concretely,
what the task is, where the boundaries are, how far it may run, who
governs it, and how we will know if tacit control is being lost.
